\subsection{搜索与剪枝}
\ojpractice{CSP202506B}{机器人复健指南：第38次，2025-06，第2题；马步八方向、最短步数与去重，第五节给解答}
\subsubsection{搜索状态和递归职责}
写递归前先明确：一次调用负责哪一步、已经作出哪些选择、接下来有哪些合法选择、什么时候记录答案。回溯是撤销本次选择，让下一条分支从同一个父状态开始。常见错误是撤销不完整，导致不同分支互相污染。

\subsubsection{例题 按字典序输出全排列}
\practice{P1706}{全排列问题}
原题要求每个数占5个字符宽度，应使用\texttt{printf("\%5d", x)}，不能直接沿用单空格输出。

输入 $1\le n\le8$，输出1到n的所有排列，每行一个，按字典序升序。
\begin{verbatim}
输入
3
输出
1 2 3
1 3 2
2 1 3
2 3 1
3 1 2
3 2 1
\end{verbatim}
path[0..depth-1] 保存已选的前缀，used[x] 标记 x 是否已在前缀中。每层从小到大尝试尚未使用的数，输出顺序自然为字典序。每个排列对应唯一一条长度 n 的路径，因此不重不漏。输出本身有 $n\cdot n!$ 个数字量级，时间 $O(n\cdot n!)$，递归与标记空间 $O(n)$。
\Code{10_permutations.cpp}

\subsubsection{剪枝必须有排除理由}
\practice{P1219}{八皇后 Checker Challenge}
逐行搜索并标记列与两条对角线；原题输出前三个方案以及方案总数。

剪枝是提前排除不可能产生合法答案或不可能改善答案的分支。可行性剪枝例如皇后问题中已占用的列与对角线；最优性剪枝例如当前代价已不小于已知最优值，且所有后续代价非负。若后续允许负代价，后一判断就可能失效。

n皇后按行放置时，记录列 c、对角线 $r+c$、另一方向 $r-c+n$ 是否已使用。每行只放一个皇后，所以行冲突由递归结构自动排除。剪枝减少搜索量，但不意味着最坏复杂度必然变成多项式。

\subsubsection{例题 迷宫最短步数}
\practice{P1443}{马的遍历}
将四方向改成马的八种跳法，按原题输出整个距离矩阵；不是只输出终点距离。

输入 $n\times m$ 字符网格，点号为可走格，井号为障碍。每步上下左右移动一格，求左上角到右下角的最少步数；不可达输出-1。$1\le n,m\le500$，起终点也可能是障碍。第一行 n、m，随后 n 行长度为 m 的字符串。
\begin{verbatim}
输入
3 4
....
.##.
....
输出
5
\end{verbatim}
\textbf{解答。}所有移动代价相等，使用 BFS。起点距离为0，邻居距离为当前距离加1；入队时就记录距离，确保同一格最多入队一次。队列按距离层次扩展，因此首次到达就是最短距离。注意这里的“步数”不包含起点，单格可走迷宫答案为0。时间与空间均为 $O(nm)$。
\Code{11_bfs.cpp}

\subsubsection{访问标记的三种含义}
\practice{P3956}{棋盘}
这道提高题包含颜色与施法规则，转移费用不全相同；要按完整状态求最小费用，不能直接提交本章普通BFS代码。

全排列中的 used 是当前路径内是否使用，返回时必须清除。普通 BFS 的 visited 表示已首次到达，通常全程不清除。记忆化的标记表示这个\textbf{完整状态的答案已算出}，应与具体到达路径无关。把这三种标记混用，会造成漏解、死循环或重复计算。

若走到同一位置时“钥匙是否拿到”影响后续，则状态必须是（位置，钥匙状态）。只用位置作为 visited 会错误合并不同状态。若边权不同，普通 BFS 不再保证最小代价；应根据权值性质考虑0--1 BFS或最短路算法，不直接套网格模板。

\subsubsection{自测与解答}
\textbf{练习1。}数组版全排列返回后，需要清空path[depth]吗？\textbf{解答。}不必。有效前缀长度由depth给出，下次选择会覆盖这一格，只有depth等于n时才输出。used[x]必须恢复为0，否则后续分支会错误地禁止使用x。若用可变长vector保存路径，则应配对push\_back与pop\_back，不能把两种实现的恢复规则混为一谈。

\textbf{练习2。}BFS为何在入队时标记，而不是出队时？\textbf{解答。}同一格可能被多个同层前驱发现，若出队才标记会重复入队。入队时记录距离既防重，也固定首次最短距离。

\textbf{练习3。}只有两个格子，边代价分别可能为1或100，按经过格子数最少就一定费用最小吗？\textbf{解答。}不一定。BFS最小化边数，只有每条边的代价相等时才等价于最小费用。选择算法前必须看目标函数。
\newpage

\subsubsection{当前团队题单训练}
\training{1112209}{搜索}
基础DFS做P1706、P1219，基础BFS做P1135、P1162、P1443；剪枝做P1120、P1731与T148457，带权状态做P1346、P3956、CF1063B。各题的题意、解答与关键代码见本篇练习题中的团队专项。
\practice{P1135}{电梯：最少按键}
\practice{P1162}{填色：从外围标记}
\practice{P1346}{电车：01BFS}
\practice{P1120}{小木棍：剪枝与恢复}
